% El zoom de un eje (mca_zoom): qué parte de la feature va al histograma y cómo
% se calcula el bin. Para el eje 1D: feature de 16 bits, AW = 13 (8192 canales),
% así que z llega a 3.
%   (a) las ventanas: en el nivel z el rango de la feature se parte en 2^z
%       ventanas iguales, alineadas, y la k-ésima ocupa el histograma entero.
%       Tres eventos: uno adentro (bin proporcional) y dos afuera, que saturan
%       en el primer y en el último canal.
%   (b) la cuenta en bits, para z = 2: los z bits altos se comparan con k, los
%       AW siguientes son el bin y el resto se descarta. Dos comparadores y un
%       mux de tres entradas, sin restas ni productos.
% Fuente: modulos/mca/rtl/mca_zoom.sv.
% Estilo: docs/figuras/estilo_tikz.tex. El zoom es de la publicación (naranja).
% El epígrafe no va en la figura: está en zoom_ventanas.txt.
%
% Compilar:   pdflatex zoom_ventanas.tex
%             pdftocairo -svg zoom_ventanas.pdf zoom_ventanas.svg
%
\documentclass[border=0pt,tikz]{standalone}
\usepackage[utf8]{inputenc}
\usepackage[T1]{fontenc}
\usepackage{lmodern}
\usepackage[none]{hyphenat}
\usepackage{amssymb}
% Estilo común de las figuras TikZ (panoramas y diagramas de bloques).
%
% Cada figura lo incluye con \input{../../figuras/estilo_tikz.tex} (las figuras
% viven en docs/<módulo>/figuras/). Un solo lugar para la paleta, los estilos de
% bloque y el formato de los parámetros: si cambia algo acá, cambia en todas.
%
% La paleta es la misma de docs/figuras/estilo_svg.py y de PALETA en
% modulos/mk/datasheet.py; si se toca un color, se toca en los tres.
%
% Requiere: \usepackage{tikz} (lo da la clase standalone con la opción tikz).

\usetikzlibrary{arrows.meta,positioning,calc,fit,backgrounds,shapes.multipart,
                decorations.markings,shadows}

% ---------------------------------------------------------------- colores ---
\definecolor{azul}{RGB}{108,172,228}      % acondicionamiento
\definecolor{verde}{RGB}{110,214,142}     % captura / histogramas en BRAM
\definecolor{amarillo}{RGB}{252,232,160}  % trigger
\definecolor{naranja}{RGB}{252,206,160}   % camino a DDR
\definecolor{grisr}{RGB}{226,226,226}     % compartido / infraestructura
\definecolor{rojo}{RGB}{240,128,128}      % bus del sistema

% ---------------------------------------------------------------- estilos ---
\tikzset{
  >={Stealth[length=2.4mm,width=1.8mm]},
  % fondo blanco siempre: el SVG se ve igual en modo claro y oscuro
  fondo/.style   = {show background rectangle, inner frame sep=8pt,
                    background rectangle/.style={fill=white, draw=none}},
  bloque/.style  = {draw=black, line width=0.6pt, rounded corners=2pt,
                    align=center, text width=2.2cm, minimum height=1.15cm,
                    inner sep=2pt, font=\scriptsize, fill=#1},
  datos/.style   = {bloque=azul!35},
  trig/.style    = {bloque=amarillo},
  captura/.style = {bloque=verde!55},
  ddr/.style     = {bloque=naranja!70},
  comp/.style    = {bloque=grisr},
  bus/.style     = {bloque=rojo!45},
  % lógica escrita dentro del módulo, sin submódulo propio
  etapa/.style   = {bloque=#1, line width=0.5pt},
  etapa/.default = verde!20,
  % lo que el bitstream no trae (atado a cero, stub, salida muda)
  mudo/.style    = {bloque=white, draw=black!45, dashed, text=black!55},
  %
  % BLOQUE CON PARÁMETROS, en el formato del datasheet: el encabezado lleva el
  % nombre y los parámetros con que lo instancia el padre (NOMBRE=valor, en
  % monoespaciada, uno por línea); el cuerpo, más claro, lo que hace.
  %   \node[datos, param={azul!35}{azul!12}] {\blk{nombre}{} \prm{A=1\\B=2}
  %                                            \nodepart{two} función};
  param/.style 2 args = {rectangle split, rectangle split parts=2,
                         rectangle split part fill={#1, #2}},
  %
  flujo/.style   = {->, line width=0.8pt},
  ctrl/.style    = {->, line width=0.6pt, black!65},
  % la línea que pasa por encima en un cruce: un halo blanco la separa
  cruce/.style   = {preaction={draw=white, line width=3.4pt, -}},
  % marca de ancho de bus: barra oblicua a mitad de la flecha, con el número
  % encima. \draw[flujo, ancho=14] ...; para otra posición, ancho en=0.3.
  ancho/.style   = {decoration={markings, mark=at position \anchopos with {
                      \draw[line width=0.6pt, -] (-0.1,-0.14) -- (0.1,0.14);
                      \node[font=\scriptsize, above=1.5pt] {#1};}},
                    postaction={decorate}},
  % la misma marca en una flecha vertical: el número va a la derecha
  anchov/.style  = {decoration={markings, mark=at position \anchopos with {
                      \draw[line width=0.6pt, -] (-0.14,-0.1) -- (0.14,0.1);
                      \node[font=\scriptsize, right=3pt] {#1};}},
                    postaction={decorate}},
  ancho en/.code = {\def\anchopos{#1}},
  sen/.style     = {font=\scriptsize\ttfamily, inner sep=1.5pt},
  nodo/.style    = {font=\tiny\ttfamily, inner sep=1pt, black!75},
  pin/.style     = {font=\scriptsize, align=center, inner sep=1pt},
  grupo/.style   = {draw=black!55, dashed, line width=0.7pt, rounded corners=4pt},
  canal/.style   = {grupo},
  contenedor/.style = {draw=black!80, line width=1.0pt, rounded corners=6pt},
  % marca numerada: cada número es una fila de la tabla del documento
  marca/.style   = {circle, fill=black, text=white, font=\scriptsize\bfseries,
                    inner sep=0pt, minimum size=4.3mm},
  %
  % DECISIONES (docs/mca/vista_general_mca.md). Cada opción es un bloque, todas
  % entran a un mux (\muxdec, abajo) y el registro que elige entra por abajo.
  % la opción que trae el reset (el comportamiento histórico): borde doble
  defecto/.style = {double=white, double distance=0.8pt, line width=0.45pt},
  % rótulo de la decisión, el mismo Dn del panorama y del documento
  decision/.style = {circle, fill=black, text=white, font=\small\bfseries,
                     inner sep=0pt, minimum size=7mm},
  %
  % LÓGICA (docs/mca/figuras/contabilidad/ y los detalles de las decisiones)
  % registro: caja gris con el triángulo de reloj en el borde de abajo
  registro/.style = {draw=black, line width=0.6pt, fill=grisr, align=center,
                     font=\scriptsize, inner sep=2pt, minimum width=#1,
                     minimum height=0.85cm,
                     path picture={\draw[line width=0.5pt]
                       ([xshift=-1.1mm]path picture bounding box.south) --
                       ([yshift=1.6mm]path picture bounding box.south) --
                       ([xshift=1.1mm]path picture bounding box.south);}},
  registro/.default = 0.8cm,
  % condición de una cadena if / else-if: "sí" sale a la derecha, "no" abajo
  pregunta/.style = {draw=black, line width=0.5pt, rounded corners=7pt,
                     fill=white, font=\scriptsize\ttfamily, align=center,
                     inner sep=2.5pt, minimum height=0.62cm, text width=#1},
  pregunta/.default = 1.9cm,
  si/.style      = {font=\tiny, inner sep=1pt, above, pos=0.12},
  no/.style      = {font=\tiny, inner sep=1pt, left, pos=0.45},
  %
  titulo/.style  = {anchor=north west, font=\small, inner sep=4pt, align=left},
  rotulo/.style  = {anchor=north west, font=\scriptsize\itshape, black!70,
                    inner sep=2pt, align=left},
}

% submódulo: nombre en monoespaciada; lógica propia: nombre en romana
\newcommand{\blk}[2]{{\ttfamily\bfseries #1}\\[1pt]#2}
\newcommand{\etq}[2]{{\bfseries #1}\\[1pt]#2}
\newcommand{\punto}[1]{\fill (#1) circle (1.3pt);}

% Parámetros en el formato del datasheet: NOMBRE=valor en monoespaciada, una
% entrada por celda, hasta cuatro columnas. Una columna: \prm{A=1\\B=2};
% dos: \prm{A=1 & B=2\\C=3}.
\newcommand{\prm}[1]{{\ttfamily\fontsize{5.6}{6.6}\selectfont
  \begin{tabular}{@{}l@{\hspace{7pt}}l@{\hspace{7pt}}l@{\hspace{7pt}}l@{}}#1\end{tabular}}}

\def\anchopos{0.5}

% Mux de una decisión: trapecio con el lado ancho (las entradas) a la izquierda.
% #1 = nombre, #2 = centro, #3 = alto en cm. Deja tres coordenadas:
%   (#1-in)  centro del lado de entradas: las flechas llegan a (#1-in |- opcion.east)
%   (#1-out) la salida
%   (#1-sel) mitad del lado de abajo: por ahí entra el registro que elige
% El rótulo de adentro es opcional: \muxdec{m}{(4,0)}{3.6} y después
% \node[...] at (m) {...}.
\newcommand{\muxdec}[3]{%
  \path #2 coordinate (#1);
  \draw[line width=0.6pt, fill=grisr]
    ($(#1)+(-0.45,{#3/2})$) -- ($(#1)+(0.45,{#3/2-0.35})$) --
    ($(#1)+(0.45,{-#3/2+0.35})$) -- ($(#1)+(-0.45,{-#3/2})$) -- cycle;
  \coordinate (#1-in)  at ($(#1)+(-0.45,0)$);
  \coordinate (#1-out) at ($(#1)+(0.45,0)$);
  \coordinate (#1-sel) at ($(#1)+(0,{-#3/2+0.175})$);}

% El mismo mux rotado 90 grados, para un flujo de arriba hacia abajo: entradas
% por el lado ancho de arriba, salida abajo y el registro que elige por la
% izquierda. #1 = nombre, #2 = centro, #3 = ancho en cm. Deja (#1-in) centro del
% lado de entradas, (#1-out) la salida y (#1-sel) mitad del lado izquierdo.
\newcommand{\muxdecv}[3]{%
  \path #2 coordinate (#1);
  \draw[line width=0.6pt, fill=grisr]
    ($(#1)+({-#3/2},0.45)$) -- ($(#1)+({#3/2},0.45)$) --
    ($(#1)+({#3/2-0.35},-0.45)$) -- ($(#1)+({-#3/2+0.35},-0.45)$) -- cycle;
  \coordinate (#1-in)  at ($(#1)+(0,0.45)$);
  \coordinate (#1-out) at ($(#1)+(0,-0.45)$);
  \coordinate (#1-sel) at ($(#1)+({-#3/2+0.175},0)$);}

% Leyenda de colores + nota. #1 = posición, #2 = color/texto, #3 = nota.
\newcommand{\leyenda}[3]{%
  \begin{scope}[shift={(#1)}, every node/.style={font=\scriptsize}]
    \coordinate (t0) at (-0.5,0);
    \foreach \c/\t [count=\i] in {#2} {
      \pgfmathtruncatemacro\j{\i-1}
      \node[draw, fill=\c, minimum width=4mm, minimum height=3mm, inner sep=0pt,
            right=5mm of t\j.east, anchor=west] (l\i) {};
      \node[right=1.5mm of l\i] (t\i) {\t};
    }
    \node[anchor=north west, black!70, text width=19cm, align=left]
      at (-0.05,-0.3) {#3};
  \end{scope}}

% Lo que dicen los parámetros, para la nota de la leyenda
\newcommand{\notaparam}{\texttt{NOMBRE=valor}: parámetros con que el padre
  instancia el bloque (se omiten los que sólo pasan el del padre sin cambio).}

\usetikzlibrary{decorations.pathreplacing}
\tikzset{
  ventana/.style = {draw=black, line width=0.5pt, fill=white, minimum height=0.42cm,
                    inner sep=0pt, font=\scriptsize},
  elegida/.style = {ventana, fill=naranja!55},
  bitc/.style    = {draw=black, line width=0.5pt, minimum width=0.62cm,
                    minimum height=0.55cm, inner sep=0pt, font=\scriptsize},
  evento/.style  = {circle, fill=black, inner sep=0pt, minimum size=1.6mm},
  rot/.style     = {font=\scriptsize, inner sep=1.5pt},
  panel/.style   = {font=\small\bfseries, anchor=north west},
}

\begin{document}
\begin{tikzpicture}[fondo]

\def\W{12}        % ancho del eje: 2^16 de la feature, y también el histograma
\def\xo{1.6}      % donde arranca el eje

% =============================================================== (a) ===
\node[panel] at (0,0.95) {(a)};

% --- el eje de la feature ---
\draw[line width=0.6pt] (\xo,0) -- (\xo+\W,0);
\foreach \f/\t in {0/$0$, 0.25/$2^{14}$, 0.5/$2^{15}$, 0.75/$3\cdot2^{14}$, 1/$2^{16}$}
  \draw[line width=0.5pt] (\xo+\f*\W,0) -- ++(0,0.1) node[above, rot] {\t};
\node[rot, anchor=east] at (\xo-0.1,0.05) {feature};

% --- las ventanas de cada nivel ---
\foreach \z/\y in {0/-0.65, 1/-1.25, 2/-1.85} {
  \node[rot, anchor=east] at (\xo-0.1,\y) {$z = \z$};
  \pgfmathtruncatemacro\nv{2^\z}
  \pgfmathsetmacro\ancho{\W/\nv}
  \foreach \k in {1,...,\nv} {
    \pgfmathtruncatemacro\kk{\k-1}
    \ifnum\z=2 \ifnum\kk=1
      \node[elegida, minimum width=\ancho cm, anchor=west] at (\xo+\kk*\ancho,\y) {$k = \kk$};
    \else
      \node[ventana, minimum width=\ancho cm, anchor=west] at (\xo+\kk*\ancho,\y) {$k = \kk$};
    \fi\else
      \node[ventana, minimum width=\ancho cm, anchor=west] at (\xo+\kk*\ancho,\y) {$k = \kk$};
    \fi
  }
}

% --- el histograma: la ventana elegida lo ocupa entero ---
\def\yh{-4.1}
\node[draw=black, line width=0.6pt, fill=naranja!25, minimum width=\W cm,
      minimum height=0.5cm, anchor=west, inner sep=0pt] (hist) at (\xo,\yh) {};
\foreach \i in {1,...,15}
  \draw[black!40, line width=0.3pt] (\xo+\i*\W/16,\yh-0.25) -- ++(0,0.5);
\node[rot, anchor=east] at (\xo-0.1,\yh) {histograma};
\node[rot, anchor=north] at (\xo,\yh-0.3) {canal 0};
\node[rot, anchor=north] at (\xo+\W,\yh-0.3) {canal 8191};

% la ventana z = 2, k = 1 se estira sobre los 8192 canales
\draw[densely dashed, black!60, line width=0.5pt] (\xo+0.25*\W,-2.06) -- (\xo,\yh+0.25);
\draw[densely dashed, black!60, line width=0.5pt] (\xo+0.5*\W,-2.06) -- (\xo+\W,\yh+0.25);

% --- tres eventos, vistos con z = 2, k = 1 ---
\def\ye{-2.55}
\node[evento] (e1) at (\xo+0.375*\W,\ye) {};
\node[evento] (e2) at (\xo+0.12*\W,\ye) {};
\node[evento] (e3) at (\xo+0.83*\W,\ye) {};
\foreach \e/\f in {e1/0.375, e2/0.12, e3/0.83}
  \draw[black!45, densely dotted, line width=0.5pt] (\xo+\f*\W,-2.06) -- (\e);
\draw[->, line width=0.7pt] (e1) -- (\xo+0.5*\W,\yh+0.25);
\draw[->, line width=0.7pt, rojo!70!black] (e2) -- (\xo+0.03,\yh+0.25);
\draw[->, line width=0.7pt, rojo!70!black] (e3) -- (\xo+\W-0.03,\yh+0.25);
\node[rot, anchor=south west] at (e1.north east) {adentro};
\node[rot, anchor=south east, rojo!70!black] at (e2.north west) {abajo: satura};
\node[rot, anchor=south west, rojo!70!black] at (e3.north east) {arriba: satura};

% =============================================================== (b) ===
\def\yb{-6.1}
\node[panel] at (0,\yb+0.95) {(b)};
\def\bw{0.62}
% bits 15..0 de la feature, de izquierda a derecha
\foreach \b in {0,...,15} {
  \pgfmathtruncatemacro\bit{15-\b}
  \ifnum\b<2
    \node[bitc, fill=grisr] (b\bit) at (\xo+\b*\bw+\bw/2,\yb) {\bit};
  \else\ifnum\b<15
    \node[bitc, fill=naranja!45] (b\bit) at (\xo+\b*\bw+\bw/2,\yb) {\bit};
  \else
    \node[bitc, fill=white, densely dashed, text=black!55] (b\bit) at (\xo+\b*\bw+\bw/2,\yb) {\bit};
  \fi\fi
}
\node[rot, anchor=east] at (\xo-0.1,\yb) {feature};
\draw[decorate, decoration={brace, amplitude=4pt, raise=1pt}, line width=0.5pt]
  (b15.north west) -- (b14.north east) node[midway, above=5pt, rot] {$z = 2$ bits altos};
\draw[decorate, decoration={brace, amplitude=4pt, raise=1pt}, line width=0.5pt]
  (b13.north west) -- (b1.north east) node[midway, above=5pt, rot] {$AW = 13$ bits: el bin};
\node[rot, anchor=south, black!55] at ($(b0.north)+(0,0.08)$) {se descarta};

% --- comparación de los bits altos con k: tres salidas ---
\coordinate (alto) at ($(b15.south)!0.5!(b14.south)$);
\node[oper] (cmp) at ($(alto)+(0,-1.25)$) {$\lessgtr$};
\node[cte, anchor=east] (kk) at ($(cmp.west)+(-0.5,0)$) {$k$};
\draw[flujo] (alto) -- (cmp.north);
\draw[flujo] (kk.east) -- (cmp.west);

% --- el mux de salida: arriba satura, adentro el bin, abajo cero ---
\def\ym{\yb-1.3}
\muxdec{m}{(\xo+12.3,\ym)}{2.3}
\node[font=\tiny, anchor=west] at ($(m-in)+(0.03,0.75)$) {$>$};
\node[font=\tiny, anchor=west] at ($(m-in)+(0.03,0)$) {$=$};
\node[font=\tiny, anchor=west] at ($(m-in)+(0.03,-0.75)$) {$<$};
\node[cte, anchor=east] (unos) at ($(m-in)+(-0.5,0.75)$) {$2^{13}-1$};
\node[cte, anchor=east] (cero) at ($(m-in)+(-0.5,-0.75)$) {$0$};
\draw[flujo] (unos.east) -- (m-in |- unos.east);
\draw[flujo] (cero.east) -- (m-in |- cero.east);
\coordinate (bsal) at ($(b13.south)!0.5!(b1.south)$);
\draw[flujo] (bsal) |- node[rot, above, pos=0.75] {bin} (m-in |- 0,\ym);
\draw[flujo] (m-out) -- ++(0.7,0) node[right, rot] {índice del bin};
\draw[ctrl] (cmp.south) -- ($(cmp.south)+(0,-0.75)$) -| (m-sel);

\end{tikzpicture}
\end{document}
